% Upper bound via path-spine + apex weight function

\subsubsection{Weight functions and prior work}

We work in the weight-function framework of Hurlbert~\cite{Hurlbert2017}, Cranston
et al.~\cite{Cranston2017}, and Yang--Yerger--Zhou~\cite{YYZ2024}; their Lemmas~1
and~2 match Section~1.2 of arXiv:2310.00580.
Fix a target~$r$.
A \emph{weight function} is $w\colon V(G)\to\mathbb{R}_{\ge0}$ with $w(r)=0$ and $w(v)>0$
for $v\ne r$.
It is \emph{valid} (rooted at~$r$) if every $r$-unsolvable configuration~$p$ satisfies
$w(p)\le w(1_G)$, where $w(p)=\sum_v p(v)w(v)$ and $w(1_G)=\sum_v w(v)$.

\begin{lemma}[Weight-function bound {\cite[ Lemma~1]{YYZ2024}}]\label{lem:wf-bound}
If $w$ is valid on $G$ rooted at~$r$ and $m=\min_{v\ne r}w(v)$, then
\[
\pi(G,r)\le\left\lfloor\frac{w(1_G)}{m}\right\rfloor+1.
\]
In particular, if $m=1$ and $w(1_G)=2^d$, then $\pi(G,r)\le 2^d+1$.
\end{lemma}

\begin{lemma}[Path weights {\cite[ Lemma~2]{YYZ2024}}]\label{lem:yyz-path}
Let $P=(v_0,v_1,\ldots,v_{k-1},r)$ be a path with $w(r)=0$ and $w(v_i)=2^i$.
For every configuration~$p$ on~$P$,
\[
\text{$p$ is $r$-solvable on $P$}\quad\Longleftrightarrow\quad w(p)\ge 2^k.
\]
\end{lemma}

\subsubsection{The path-spine weight function}\label{sec:path-spine-wf}

Fix $G=T(3,s_1,s_2)$, $d=s_1+s_2+1$, and target $r=u_R$.
Let
\[
P=(u_L,w_1,\ldots,w_{s_1-1},a,b,z_1,\ldots,z_{s_2-1},u_R)
\]
be the Hamiltonian path of length~$d$ (so $|V(P)|=d+1$).
Define $w_P$ by
\[
w_P(u_R)=0,\qquad w_P(v)=2^i\ \text{when $v$ is the $(i+1)$-st vertex of $P$ starting from $u_L$},
\]
and $w_P(c)=0$.
Set $w_c(c)=1$ and $w_c(v)=0$ for $v\ne c$, and define
\[
w=w_P+w_c.
\]

\begin{lemma}[Normalization]\label{lem:w-normalization}
For $w$ as above, $w(1_G)=2^d$ and $\min_{v\ne u_R}w(v)=1$.
\end{lemma}

\begin{proof}
The path contributes $1+2+\cdots+2^{d-1}=2^d-1$, and $w_c(c)=1$.
\end{proof}

\subsubsection{Validity}\label{sec:apex-feed}

Set $Q=2^d$.
For a configuration~$p$ on~$G$, let $w_P(p)$ and $w_c(p)=p(c)$ denote the path and
apex contributions.

\begin{lemma}[Path case]\label{lem:path-case}
If $w_P(p)\ge Q$, then $p$ is $u_R$-solvable on~$G$.
\end{lemma}

\begin{proof}
Apply Lemma~\ref{lem:yyz-path} on the path subgraph~$P$.
Every pebbling move on~$P$ is also legal on~$G$.
\end{proof}

\begin{lemma}[Apex feeding]\label{lem:apex-feed}
Suppose $w_P(p)<Q$ but $w(p)>Q$.
Then $p$ is $u_R$-solvable on~$G$.
\end{lemma}

\begin{proof}
Since $w(p)=w_P(p)+p(c)>Q$ and $w_P(p)\le Q-1$, we have $p(c)\ge 2$.
Let $\delta=Q-w_P(p)\ge 1$.

Vertex~$c$ is adjacent only to $a$ and~$b$, both on~$P$.
A \emph{feeding move} is $c\to a$ or $c\to b$; it is legal exactly when $p(c)\ge 2$.
Each feeding move removes two pebbles from~$c$ and adds one to a spine vertex of~$P$,
so $w_P$ increases by $2^{s_1}$ if the move is $c\to a$ and by $2^{s_1+1}$ if it is
$c\to b$; in either case $\Delta w_P\ge 2^{s_1}$.

Perform feeding moves repeatedly, in any order, while the current configuration
has at least two pebbles on~$c$.
This defines a legal pebbling sequence: at step~$j$, before the $j$-th feeding move
the configuration has $p_j(c)\ge 2$, so the move is permitted, and afterward
$p_j(c)$ decreases by exactly~$2$.
After $k=\lfloor p(c)/2\rfloor$ such moves we reach a configuration~$p'$ with
$p'(c)\in\{0,1\}$ and
\[
w_P(p')\ge w_P(p)+k\cdot 2^{s_1}.
\]

Because $p(c)\ge \delta+1$ with $\delta=Q-w_P(p)\ge1$, a direct check shows
$w_P(p')\ge w_P(p)+\lfloor p(c)/2\rfloor\cdot 2^{s_1}\ge Q$.
Lemma~\ref{lem:path-case} applies to~$p'$, so $p'$ is $u_R$-solvable; hence so is~$p$.
\end{proof}

\begin{lemma}[Validity of $w$]\label{lem:wf-valid}
The weight function $w$ from Section~\ref{sec:path-spine-wf} is valid on
$T(3,s_1,s_2)$ rooted at~$u_R$.
\end{lemma}

\begin{proof}
Let $p$ be $u_R$-unsolvable.
If $w_P(p)\ge Q$, Lemma~\ref{lem:path-case} gives a contradiction.
If $w_P(p)<Q$ but $w(p)>Q$, Lemma~\ref{lem:apex-feed} gives a contradiction.
Therefore $w(p)\le Q=w(1_G)$.
\end{proof}

\subsubsection{Main upper-bound and exact formula}

\begin{theorem}[Upper bound]\label{thm:upper-bound}
For all $s_1,s_2\ge 1$ with $d=s_1+s_2+1$,
\[
\pi(T(3,s_1,s_2))\le 2^d+1.
\]
\end{theorem}

\begin{proof}
By Lemmas~\ref{lem:wf-valid} and~\ref{lem:w-normalization},
Lemma~\ref{lem:wf-bound} gives $\pi(G,u_R)\le 2^d+1$.
By symmetry (reflect the graph, swapping the left and right tails), the same
construction with target~$u_L$ gives $\pi(G,u_L)\le 2^d+1$.
Lemma~\ref{lem:leaf} then yields $\pi(G)\le 2^d+1$.
\end{proof}

\begin{theorem}[Exact pebbling number]\label{thm:exact}\label{thm:main}
For all $s_1,s_2\ge 1$ with $d=s_1+s_2+1$,
\[
\pi(T(3,s_1,s_2))=2^d+1=2^{\operatorname{diam}(T(3,s_1,s_2))}+1.
\]
\end{theorem}

\begin{proof}
Theorem~\ref{thm:lower-bound} gives $\pi(T(3,s_1,s_2))\ge 2^d+1$, and
Theorem~\ref{thm:upper-bound} gives $\pi(T(3,s_1,s_2))\le 2^d+1$.
The lower and upper bounds coincide, hence
$\pi(T(3,s_1,s_2))=2^d+1$.
Since $d=\operatorname{diam}(T(3,s_1,s_2))$ within this family, the formula depends on
$d$ alone.
\end{proof}

Configuration~$D^*$ has $w(D^*)=2^d$ at $S=2^d$, so the weight bound is tight.
